% Définition du répertoire contenant les images
\graphicspath{{IMAGE/}}


% FRAME Intro
\begin{frame}

\includegraphics[width=12cm]{Logos.pdf}

\vfill

\begin{center}

\vspace*{1.5cm}

\LARGE
\textbf{Accessibility}

\vspace*{1.5cm}
 DELHI GIS-R School


\large
9-12\up{th} April 2019

\vspace*{1.5cm}


\textbf{Hadrien Commenges \& Paul Chapron}

{\small

\vspace*{0.1cm}

\url{hadrien.commenges@univ-paris1.fr}

\url{paul.chapron@ign.fr}
}

\end{center}

\end{frame}


% FRAME
\begin{frame}{Use case}


\begin{block}{Accessibility}
Accessibility depicts the potential access of an agent to a resource.

\end{block}

\textbf{two meanings:}

\begin{enumerate}
  \item \textbf{Access to a location:} required effort for agents to get to a location ($\approx$ service area).Such a measure depends on the network (infrastructure + services) and the agent's capacity to travel.
  \item \textbf{Access to a resource:} required effort for agents to get to a resource. Suche a measure depends on the network, the agent's capacity to travel and the resource \textbf{spatial distribution}.
\end{enumerate}

~

~

\textbf{What kind of geographical information is concerned ?}

$\rightarrow$ \textit{TYPE 1 - Geographical Objects}

$\rightarrow$ \textit{TYPE 2 - Occurrences}

$\rightarrow$ \textit{TYPE 3 - Measure points}

$\rightarrow$ \textit{TYPE 4 - Statistics}

$\rightarrow$ \textit{TYPE 5 - Interaction Measures}


\end{frame}



% FRAME
\begin{frame}{Distance and interaction}

\begin{itemize}
  \item \textbf{friction} (\textit{impedance}) is the key-concept, distance is a measure of friction.
  \item Spatial interaction is adressed via \textbf{interactions between locations} (flows) or \textbf{attraction/influence of a location on the others}. 
  \item Accessibility models rely on the latter.
\end{itemize}

\end{frame}

